\section{Concluding remarks}

Our main goal was to answer the question of Montgomery and Shanbrom about integrability of~the (simple) Kepler problem on the Heisenberg group. The answer turned out to be negative, but several generalisations became immediately apparent. First, the potential had a specific radial/axial symmetry, and a whole general class of such potentials could be included; second, and more important, the two-body problem could be formulated in a natural way. We thus extended the analysis, and managed to show, that with reasonable assumptions those extensions were also non-integrable.

We note that potentials not satisfying condition \eqref{eq:con} can be found, such as
\begin{equation*}
 V = z^2\big(x^2+y^2\big)^2 = z^2\big(\rho^2-16z^2\big), \qquad\text{or}\qquad V=\big(x^2+y^2\big)^2R(z,\rho),
\end{equation*}
where $R(z,\rho)$ is not divisible by $\big(x^2+y^2\big)^2$. Integrability of these potentials remains an open question.
One possible way of investigation of such cases is the application of a variant of the direct method. However, we were unable to find any integrals which were polynomials of low degree in momenta. It remains an open question whether our result can be extended to a wider functional class of first integrals, but each case requires a completely different set of methods than those used here, and as such is a subject for separate investigation.
